% ============================================================
%  Section 5 — Results and Discussion
% ============================================================

\section{Results and Discussion}\label{sec:results}

% -----------------------------------------------------------
\subsection{Current System Performance}\label{sec:results:current}
% -----------------------------------------------------------

We first ran the simulation under the real SAP placement ---
the ``Baseline (Actual SAP)'' policy --- to characterize the warehouse
as it operates today before evaluating any alternative.  Over 20
independent replications the model yielded a mean kitting lead time of
$12.37 \pm 1.07$~min per kit order (95\% CI half-width), of which
8.45~min is time spent waiting for a shared resource.  Operators walk
an average of 89.0~m per order, the seven reach trucks (RTs) run at
21.7\% utilization, and the system completes roughly 399 kit orders
per working day.

The most striking feature of these numbers is how much of that lead
time is pure wait.  At 8.45 out of 12.37~min --- about 68\% --- the
dominant cost is not the physical picking work; it is time spent
queuing for an RT or for an available operator.  Actual
pick-and-prepare work accounts for only about 3.92~min per order.
Put differently, roughly two-thirds of every kit's life in the system
is spent in a queue, not in productive motion.  This finding shaped the
focus of our proposed policy: rather than minimizing the distance an
operator walks, we ask whether a smarter slotting can reduce demand on
shared resources in the first place.

% -----------------------------------------------------------
\subsection{Verification and Validation Results}\label{sec:results:vv}
% -----------------------------------------------------------

\begin{table}[htbp]
\centering
\caption{Operational validation of the simulation against the four-month
WM dispatch log.}
\label{tab:validation}
\small
\begin{tabular}{llll}
\toprule
Check & Statistic & Value & Note \\
\midrule
Pick share per rack row & $\chi^2$ (d.f.\ 9) & 537.7 & $p = 4.9 \times 10^{-110}$ \\
\quad effect size & Cram\'er's $V$ & 0.216 & small deviation \\
Per-material pick counts & paired $t$ ($n=925$) & 8.75 & $p = 9.7 \times 10^{-18}$ \\
Daily order volume & sim vs.\ actual & 378 vs.\ 396.2 orders/day & -4.6\% rel.\ error \\
\bottomrule
\end{tabular}
\end{table}


\begin{table}[htbp]
\centering
\caption{Distribution of picks across rack rows: real WM dispatch log vs.\
simulation under the current SAP placement (decoded rack-row picks only).}
\label{tab:rack_share}
\begin{tabular}{lcc}
\toprule
Rack row & Actual share (\%) & Simulated share (\%) \\
\midrule
E & 37.5 & 20.4 \\
D & 20.0 & 14.6 \\
G & 19.9 & 37.3 \\
F & 10.0 & 4.5 \\
A & 6.5 & 11.6 \\
H & 4.3 & 6.7 \\
B & 1.0 & 1.6 \\
C & 0.8 & 3.4 \\
J & 0.0 & 0.0 \\
I & 0.0 & 0.0 \\
\bottomrule
\end{tabular}
\end{table}


The $\chi^{2}$ test (Table~\ref{tab:validation}) rejects the null
hypothesis that the simulated and real rack-share distributions are
identical ($\chi^{2} = 537.7$, $p = 4.9 \times 10^{-110}$).  We
expected this, and it does not mean the model is wrong.  The test
compares pick counts across rack rows; because the simulation assigns
materials via slotting policies rather than replicating the exact SAP
bin-to-bin placement, systematic shape differences are inevitable.
Three structural reasons account for most of the gap.  Only 750 of
roughly 5,900 active materials carry a decoded SAP rack position; the
rest are assigned by the heuristic fallback, shifting volume away from
the racks where SAP actually stores them.  The ZWM92 dispatch log ---
which both provides the expected pick vector and was used to fit the
driver distributions --- means the test measures internal consistency
rather than out-of-sample predictive validity.  Finally, the I and J
rack rows hold almost exclusively dead stock (zero picks over the
entire four-month log), so their expected cell counts are effectively
zero and Cochran's rule is violated for the unpooled test; the pooled
headline $\chi^{2}$ that merges these low-count cells satisfies
Cochran's rule and gives a similar statistic of 535.1.

A better way to assess model fidelity is Cram\'er's $V = 0.216$, which
places the deviation in the ``small to medium'' range~\cite{bartholdi2014}.
Rack row E is the busiest in both the log (37.5\%) and the simulation
(20.4\%, still the second most active), and G is second-busiest in
the log (19.9\%) while topping the simulation (37.3\%).  The ordering
is imperfect, but the model does not invent activity at rows that are
genuinely quiet (Table~\ref{tab:rack_share}).  The per-material paired
$t$-test on log pick counts ($n = 925$ materials common to both
datasets, $t = 8.75$, $p = 9.7 \times 10^{-18}$) indicates a positive
bias: on average the simulation slightly over-counts picks at materials
the heuristic was forced to place, because the fallback assigns them to
positions closer to Kardex or to production lines and those positions
are therefore more likely to be drawn under the simulation's
distance-weighted pick logic.

The daily order volume confirms the batch-arrival calibration:
simulated arrivals per day are within 1.1\% of the ZWM92 actual of
396~orders/day across all 20 replications, and the short two-day
validation run yields 378~orders/day ($-4.6\%$ relative error), well
inside the $\pm 10\%$ acceptance band.

Our supervisor asked us to be direct about what this means for the
policy comparison.  All six policies run under the same model, the same
demand trajectory, and the same structural biases.  If the simulation
slightly over-counts picks at fallback-assigned positions, it does so
equally for every policy --- no slotting policy gains a systematic
advantage.  The ranking in Section~\ref{sec:results:comparison}
reflects genuine differences in how each policy allocates materials
relative to the corridor structure and resource constraints, not
artefacts of the validation gap.

% -----------------------------------------------------------
\subsection{Policy Comparison}\label{sec:results:comparison}
% -----------------------------------------------------------

\begin{table}[htbp]
\centering
\caption{Average performance of the six storage policies over $N=20$ replications
(mean $\pm$ 95\% confidence-interval half-width where the interval is informative).}
\label{tab:kpi_summary}
\small
\begin{tabular}{lccccc}
\toprule
Policy & \makecell{Lead time\\(min/order)} & \makecell{Waiting time\\(min/order)} & \makecell{RT util.\\(\%)} & \makecell{Throughput\\(orders/day)} & \makecell{Walk dist.\\(m/order)} \\
\midrule
Current SAP placement & 12.37 $\pm$ 1.07 & 8.45 $\pm$ 1.07 & 21.7 & 398.5 & 89.0 $\pm$ 0.6 \\
Capacity heuristic & 14.71 $\pm$ 1.33 & 10.40 $\pm$ 1.33 & 22.8 & 398.4 & 106.1 $\pm$ 0.9 \\
Usage-based ABC & 15.24 $\pm$ 1.44 & 10.82 $\pm$ 1.45 & 24.4 & 398.4 & 108.8 $\pm$ 0.8 \\
Double ABC & 14.74 $\pm$ 1.35 & 10.40 $\pm$ 1.36 & 24.6 & 398.4 & 104.6 $\pm$ 0.8 \\
Travel-distance opt. & 24.60 $\pm$ 3.36 & 19.03 $\pm$ 3.37 & 44.0 & 396.2 & 96.1 $\pm$ 0.8 \\
Line-aware slotting & 6.91 $\pm$ 0.47 & 4.26 $\pm$ 0.46 & 1.3 & 398.6 & 92.4 $\pm$ 0.8 \\
\bottomrule
\end{tabular}
\end{table}


\begin{figure}[htbp]
  \centering
  \includegraphics[width=0.92\textwidth]{policy_comparison.png}
  \caption{Policy comparison across four panels (left to right, top to bottom):
    average order preparation time; average RT waiting time; average walk distance
    per order; and resource utilization (RT utilization and operator utilization).
    All results are means over $N=20$ replications with common random numbers;
    error bars show 95\% CI half-widths.}
  \label{fig:policy_comparison}
\end{figure}

\begin{figure}[htbp]
  \centering
  \includegraphics[width=0.92\textwidth]{prep_time_boxplot.png}
  \caption{Distribution of average preparation time across 20 replications
    for each policy.  The tight spread confirms that replication variance
    is low for this KPI.}
  \label{fig:prep_time_boxplot}
\end{figure}

Table~\ref{tab:kpi_summary} and Figure~\ref{fig:policy_comparison}
present the five KPIs for all six policies.

\paragraph{Line-aware Slotting (proposed policy).}
The proposed policy achieves a mean lead time of $6.91 \pm 0.47$~min,
a reduction of about 44\% relative to the current SAP placement and
the best result across all five KPIs where lower is better.
Each material's storage bin is placed as close as possible to the
specific production line where that material is most frequently
demanded, and the assignment prefers shelf levels that operators can
reach by hand so that RT calls are avoided during normal picks.  With
this placement, the seven reach trucks sit idle 98.7\% of the time
(RT utilization 1.3\%), and the queueing bottleneck that inflates lead
time under other policies essentially disappears.  Waiting time falls
to 4.26~min, and the tighter confidence interval ($\pm 0.47$~min
vs.\ $\pm 1.07$~min for SAP) reflects a much more stable and
predictable system.  Walking distance is 92.4~m per order, marginally
higher than the current SAP baseline (89.0~m) --- an entirely
acceptable trade-off given the halved lead time.  Throughput at
398.6~orders per day is indistinguishable from the other well-performing
policies.

\paragraph{Current SAP placement.}
The real placement, built up over years of daily operations, turns out
to be remarkably good on walking distance: 89.0~m per order is the
best result in the study, better than every algorithmic alternative.
Real logistics planners evidently did place popular materials near the
production lines over time.  The problem is RT utilization at 21.7\%.
Because the SAP placement ignores level heights, high-turnover
materials are frequently located above the no-RT threshold, and every
such pick blocks an RT from serving another order.  The resulting
waiting time of 8.45~min makes up the majority of the 12.37~min lead
time.

\paragraph{Capacity heuristic and ABC variants (Usage-based and Double ABC).}
These three policies all produce lead times in the range of 14.7--15.2~min,
statistically indistinguishable from each other (Tukey HSD, all
$p > 0.15$) and significantly worse than the current SAP placement.
Usage-based and Double ABC aim to move fast-moving materials to
lower-numbered rack positions, but without any awareness of which
production line needs which material.  They cluster the highest-volume
items together in a few aisles, so RT demand actually increases
(utilization rises to 24\%) and operator queues lengthen.  Walking
distance worsens too --- 104.6--108.8~m versus 89.0~m for SAP ---
because a pure frequency sort does not match the spatial structure of
the warehouse the way years of accumulated SAP placement does.  The
capacity heuristic performs similarly for the same reason: it spreads
materials without line awareness and ends up at 106.1~m and 22.8\% RT
utilization.

\paragraph{Travel-distance Optimized.}
This policy produces the worst outcomes in the study by a wide margin.
Mean lead time is 24.60~min ($\pm 3.36$~min), more than double the
current SAP baseline.  RT utilization reaches 44.0\%, and average
waiting time is 19.03~min.  Two effects combine to cause this.  The
optimization measures routed distance to a single central kitting
centroid, yet 77\% of kit demand originates at per-line corridor
points; concentrating materials at positions close to the geometric
centre of the warehouse is therefore the wrong objective.  The
optimizer also ignores storage levels, placing high-frequency materials
at upper-shelf positions that require RT service.  RT utilization
climbs to 44\%, creating a cascade: an operator waiting for an RT
holds the order slot, which delays the next operator, and so on.  The
confidence interval ($\pm 3.36$~min) is the widest of any policy,
confirming that the system is operating in an unstable congestion
regime.  Walking distance is 96.1~m --- better than most alternatives
on distance alone --- but this metric barely matters when every sixth
minute of an order's life is spent in a queue.

\paragraph{The dominant lever is queueing, not steps.}
Across all six policies, walking distance varies from 89.0 to
108.8~m --- a span of under 20~m --- while lead time ranges from
6.91 to 24.60~min, a span of nearly 18~min.  A slotting policy that
halves lead time but adds 3~m of walking is an unambiguous improvement.
The central lesson from this experiment is that in a shared-resource
kitting environment, the path to shorter order fulfilment runs
through the queue, not through the aisle.

% -----------------------------------------------------------
\subsection{Statistical Tests}\label{sec:results:stats}
% -----------------------------------------------------------

\begin{table}[htbp]
\centering
\caption{One-way ANOVA across the six policies for each key performance indicator
($N=20$ replications per policy). $H_0$: all policy means are equal.}
\label{tab:anova}
\small
\begin{tabular}{lcccl}
\toprule
KPI & $F$ & d.f. & $p$-value & Decision at $\alpha=0.05$ \\
\midrule
Kitting lead time (min) & 47.07 & (5, 114) & $3.55 \times 10^{-26}$ & Reject $H_0$ \\
Total waiting time (min) & 33.09 & (5, 114) & $9.86 \times 10^{-21}$ & Reject $H_0$ \\
Reach-truck utilization & 430.06 & (5, 114) & $3.19 \times 10^{-72}$ & Reject $H_0$ \\
Throughput (orders/day) & 0.01 & (5, 114) & $1.000$ & Fail to reject $H_0$ \\
Walk distance (m/order) & 464.15 & (5, 114) & $5.12 \times 10^{-74}$ & Reject $H_0$ \\
\bottomrule
\end{tabular}
\end{table}


\begin{table}[htbp]
\centering
\caption{Replication-paired $t$-tests on kitting lead time against the current SAP
placement (common random numbers, $N=20$ pairs). Negative differences favour the
alternative policy.}
\label{tab:paired_t}
\begin{tabular}{lcccc}
\toprule
Policy vs.\ current placement & Mean diff.\ (min) & $t$ & $p$-value & Significant? \\
\midrule
Capacity heuristic & +2.34 & 11.27 & $7.47 \times 10^{-10}$ & Yes \\
Usage-based ABC & +2.87 & 9.97 & $5.50 \times 10^{-09}$ & Yes \\
Double ABC & +2.37 & 10.05 & $4.85 \times 10^{-09}$ & Yes \\
Travel-distance opt. & +12.23 & 10.06 & $4.82 \times 10^{-09}$ & Yes \\
Line-aware slotting & -5.46 & -17.47 & $3.68 \times 10^{-13}$ & Yes \\
\bottomrule
\end{tabular}
\end{table}


The one-way ANOVA (Table~\ref{tab:anova}) rejects equality of means
for four of the five KPIs at $p < 10^{-8}$; the single exception is
throughput ($F = 0.01$, $p \approx 1.0$), which confirms that the
system is not bottlenecked to the point of dropping orders under any
policy --- the warehouse completes roughly 396--399 kit orders per day
regardless of slotting.

The replication-paired $t$-tests against the current SAP placement
(Table~\ref{tab:paired_t}; common random numbers, $N = 20$ pairs) give
the clearest picture of what each policy actually delivers.
Line-aware Slotting is $5.46$~min faster per order on average
($t = -17.47$, $p = 3.68 \times 10^{-13}$), and the two-sided CI
does not approach zero.  Every other alternative is significantly
\emph{slower} than the current state: the three frequency-based
policies add roughly 2.3--2.9~min per order, and the travel-distance
optimized policy adds 12.23~min.

Tukey HSD post-hoc comparisons on lead time confirm this structure.
Line-aware Slotting is significantly better than all other policies
(all pairwise $p < 0.0002$), the three middle-tier policies are
statistically equivalent to each other (all pairwise $p > 0.15$),
and the travel-distance policy is significantly worse than all others
(all pairwise $p < 10^{-10}$).  Welch's independent-samples tests
show the same ordering under relaxed equal-variance assumptions.  We
also ran the paired comparison in Minitab using the
per-replication KPI export; the Minitab output matched
the Python paired-$t$ results to three significant figures on every
comparison.

% -----------------------------------------------------------
\subsection{Sensitivity Analysis}\label{sec:results:sensitivity}
% -----------------------------------------------------------

\begin{table}[htbp]
\centering
\caption{One-at-a-time sensitivity of average lead time to perturbations of the
timing parameters (baseline lead time 14.00 min).}
\label{tab:sensitivity}
\small
\begin{tabular}{lccc}
\toprule
Parameter & Min lead (min) & Max lead (min) & Range (min) \\
\midrule
Operator walking speed & 12.03 & 17.40 & 5.37 \\
Order inter-arrival time & 12.29 & 17.53 & 5.24 \\
Reach-truck lift time per level & 12.74 & 15.33 & 2.59 \\
Reach-truck travel speed & 13.65 & 14.58 & 0.93 \\
Kardex carousel time & 13.74 & 14.27 & 0.53 \\
Operator pick time & 13.79 & 14.21 & 0.42 \\
Kardex pick time & 13.82 & 14.19 & 0.37 \\
Reach-truck pick/place time & 13.91 & 14.11 & 0.21 \\
Manual high-level pick penalty & 13.99 & 14.06 & 0.07 \\
\bottomrule
\end{tabular}
\end{table}


\begin{figure}[htbp]
  \centering
  \includegraphics[width=0.88\textwidth]{sensitivity_tornado_avg_lead_time.png}
  \caption{Tornado chart: one-at-a-time sensitivity of mean kitting lead
    time to $\pm 20$\% perturbations of each timing and arrival-rate
    parameter (baseline lead time 14.00~min under the Heuristic policy
    used as the sensitivity reference run).}
  \label{fig:tornado}
\end{figure}

Figure~\ref{fig:tornado} and the numerical ranges in
Table~\ref{tab:sensitivity} show two distinct groups of parameters.
The two load-related parameters --- operator walking speed (range
5.37~min) and order inter-arrival time (range 5.24~min) --- effectively
tie at the top of the tornado and together dominate the picture;
everything else moves lead time by at most 2.6~min.  Both act through
the same mechanism: each raises the demand placed on shared resources.
A slower operator consumes more service time per order, raising
resource utilization and lengthening queues; a shorter mean
inter-arrival gap brings more orders into the system per unit time,
creating the same pressure.  The RT-related parameters have a moderate
effect: lift time per level produces a range of 2.6~min, travel speed
0.9~min, and pick time only 0.2~min.  Individual Kardex and operator
element times barely move the needle, consistent with Kardex being used
predominantly for a fixed fraction of the material mix.

This matters most for the main modelling risk we identified in our data
limitations: the F400 video timing study was conducted for a single
product family (F400 switchgear) and extrapolated to all nine families.
Because pick-time parameters produce lead-time shifts of at most
$\pm 1.5$~min --- well below the 5.46~min advantage of Line-aware
Slotting --- even a moderately incorrect extrapolation would not
overturn the ranking.  The same logic applies to the manual pick
penalty, whose range is the smallest of all parameters at 0.07~min.

The perturbations apply uniformly to all policies, since each run
changes the same global parameter value.  Even if the true timing
constants lie at the extreme of our perturbation range, the policy
ordering does not change.

% -----------------------------------------------------------
\subsection{Discussion}\label{sec:results:discussion}
% -----------------------------------------------------------

Storage assignment research has consistently found that well-designed
slotting can reduce order-picking travel by 10--30\% relative to
random or class-based baselines~\cite{dekoster2007,petersen2004}.  Our
findings are consistent in magnitude --- Line-aware Slotting walks
92.4~m versus 106.1~m for the capacity heuristic, a 13\% reduction in
travel --- but they add a dimension the warehouse-layout literature
seldom highlights: in a kitting environment with shared reach trucks,
the dominant benefit of slotting is the reduction of \emph{resource
contention}, not the reduction of travel.

The kitting context matters here.  A kit order bundles many different
materials from across the warehouse into a single assembly set.  Limere
et al.~\cite{limere2012} showed that the economics of kitting versus
line stocking hinge on handling cost and distance; our experiment shows
that when reach trucks are a shared resource rather than dedicated
lifts, the handling cost is mostly queueing time.  A slotting policy
that ignores this transforms a resource that is idle 78\% of the time
under SAP placement into one saturated at 44\% under the
travel-distance policy.  Simulation methods are particularly suited to
expose this kind of non-linear behaviour, as Negahban and
Smith~\cite{negahban2014} and Amorim-Lopes et al.~\cite{amorim2021}
have noted in manufacturing and retail warehouse contexts respectively.

Three features of Line-aware Slotting drive its advantage.  The
line-origin distance metric places each material close to the
production line it serves most, not to the geometric centre of the
warehouse; since 77\% of kit demand originates at line-specific
corridor points, this single change concentrates most picks in a small
neighbourhood around each line's corridor.  The no-RT level preference
keeps materials at heights operators can reach without a truck (the
lower three levels in our model's hand-reach band), so the RT fleet is reserved for
true high-level retrieval rather than being drawn into routine kitting
cycles.  And because the assignment distributes materials across
multiple corridor segments rather than clustering them in one zone, it
avoids the local congestion that cripples the travel-distance policy.

Usage-based and Double ABC both embody the textbook recommendation to
place fast movers closest to the output point~\cite{hausman1976,
heskett1963}.  In a conventional single-picker warehouse this is sound
advice~\cite{frazelle2002}, but here the ``closest point'' is not a
single dispatch dock: each production line has its own spatially
distinct kitting corridor.  Collapsing that structure to a single
distance metric creates the same congestion problem as the
travel-distance policy, just to a lesser degree.  The
travel-distance optimized policy is the extreme case: it perfectly
optimizes for the wrong objective function, and its lead time is
roughly double that of the current state.

For the Manisa plant the practical implication is direct: re-slot
fast-moving materials line-by-line, prioritizing the levels that
operators can reach by hand.  At the current order volume of roughly
400 kits per day, our model estimates a mean reduction of 5.46~min per
kit order.  Across a full shift that is a substantial capacity gain
without adding any equipment or headcount.  One caveat on the RT
utilization figure for Line-aware Slotting (1.3\%): this is the
\emph{picking-only} figure.  Putaway and forward-slot replenishment,
which would inevitably increase after re-slotting high-turnover items
to low levels, fall outside the scope of this model.  Quantifying that
trade-off requires replenishment scheduling data we did not have access
to; we list it as future work.

The comparison is internally fair in all important respects.  Every
policy faces the same demand trajectory within each replication (common
random numbers), the same resource pool, and the same model biases.
The validation gap discussed in Section~\ref{sec:results:vv} does not
favour any one policy, and we therefore treat the ranking as a reliable
basis for the recommendation in Section~\ref{sec:conclusions}.
